\chapter{浮点数、数值线性代数与病态问题}
\label{ch:numerical-stability}

数值计算的危险之处在于，程序往往会正常结束并返回一个看似精确的数字。本章从机器实际存储的浮点数出发，区分问题的条件性、算法的稳定性和实现的舍入误差，再通过消去、求和、线性系统和最小二乘的例子说明三者如何叠加。读者最终要学会问：这个结果是原问题的答案，还是某个邻近问题的稳定答案？

所有数值建议都以可解释的尺度和独立基准为前提。更高精度、更多小数位或更宽容差都不能替代对算法和问题本身的诊断。

\section{表示、问题与算法}

binary64 把有限数编码为符号、11 位指数和 52 位显式尾数；正常数还含一个隐含的首位，因此有效精度为 53 位。相邻可表示数之间的距离随指数变化。若 $x$ 是正常数，$\operatorname{ulp}(x)$ 表示其末位单位；机器精度
\[
  \varepsilon_{\mathrm{mach}}=\operatorname{nextafter}(1,+\infty)-1=2^{-52},
\]
而舍入到最近数时常用的单位舍入误差是 $u=\varepsilon_{\mathrm{mach}}/2=2^{-53}$。两者差一倍，报告时必须说明定义。证据环境中，机器精度为 $2.220\times10^{-16}$，单位舍入误差为 $1.110\times10^{-16}$。

对结果仍处于正常范围、没有上溢或下溢的基本运算，常用模型是
\[
  \operatorname{fl}(a\circ b)=(a\circ b)(1+\delta),\qquad |\delta|\le u.
\]
它是单次运算模型，不等于“任意长程序只有一个 $u$ 的误差”。次正规数填补零与最小正常数之间的间隙，但相对误差界会失效；再小则逐渐下溢为零。$+0$ 与 $-0$ 比较相等却能保留方向信息；无穷可表达有限数溢出或除零；NaN 与任何数（包括自身）的相等比较均为假。默认舍入模式通常是“最近、偶数优先”，向上、向下和朝零舍入会改变误差方向。

十进制小数通常不能由有限二进制精确表示。金额不能因此一律改用容差：协议字段、计数和以最小货币单位记录的整数仍应精确比较；收益、回归系数等近似量才需要带单位的误差策略。

\MFQLead{一个浮点数究竟表示了什么}

以十进制 $0.1$ 为例，二进制展开是无限循环的
\[
  0.1=(0.0001100110011\ldots)_2.
\]
binary64 只能截取有限位并舍入，所以 $\operatorname{fl}(0.1)\ne0.1$；但误差小于一个单位舍入误差乘以 $0.1$。这解释了为什么
\[
  (0.1+0.1+0.1)-0.3
\]
可能得到一个很小但非零的数，而不是说明加法定律失效。若协议要求金额精确，应先把金额转成整数分；若对象本来就是估计量，则应报告绝对/相对误差的尺度，而不是把所有比较改成字符串相等。

相邻浮点数的距离可以用 $\operatorname{ulp}$ 读出：在 $[1,2)$ 内相邻数间隔为 $2^{-52}$，在 $[2,4)$ 内间隔变为 $2^{-51}$。一个适合 $10^{-8}$ 的阈值可能完全不适合 $10^8$。

\MFQLead{三个对象：问题、算法与实现结果}

设精确问题为 $y=f(x)$，输入扰动为 $\Delta x$，计算结果为 $\widehat y$。绝对前向误差是 $\|\widehat y-y\|$，相对前向误差是
\[
  \eta_{\mathrm f}=\frac{\|\widehat y-y\|}{\|y\|}.
\]
后向误差寻找最小的 $\Delta x$，使 $\widehat y=f(x+\Delta x)$；它问“计算结果精确解了哪个邻近问题”。问题的相对条件数
\[
  \kappa_f(x)=\lim_{\epsilon\to0}\sup_{\|\Delta x\|\le\epsilon\|x\|}
  \frac{\|f(x+\Delta x)-f(x)\|/\|f(x)\|}{\|\Delta x\|/\|x\|}
\]
衡量精确映射本身的敏感度。后向稳定算法只保证后向误差小；一阶上通常有
\[
  \text{相对前向误差}\lesssim \kappa_f(x)\times\text{相对后向误差}.
\]
因此稳定算法无法修复病态问题，良态问题也可能被不稳定公式算坏。误差界还必须写清范数、尺度与单位；当真值接近零时，相对误差可能没有解释力。

\begin{proposition}[线性系统的条件数放大]
若 $Ax=b$ 可逆，右端扰动为 $\Delta b$，则
\[
  \frac{\|\Delta x\|}{\|x\|}
  \le \kappa(A)\frac{\|\Delta b\|}{\|b\|}
  \quad\text{（在相容范数下）},
\]
其中 $\Delta x=A^{-1}\Delta b$。等号通常只有在扰动对准最坏奇异向量方向时才近似达到；随机的小扰动可能远小于条件数上界。
\end{proposition}
\begin{proof}
由 $\Delta x=A^{-1}\Delta b$ 和 $b=Ax$，
\[
 \frac{\|\Delta x\|}{\|x\|}
 \le\|A^{-1}\|\frac{\|\Delta b\|}{\|x\|}
 \le\|A^{-1}\|\|A\|\frac{\|\Delta b\|}{\|b\|}.
\]
在二范数下，$\|A\|\|A^{-1}\|$ 就是最大与最小奇异值之比。这个界描述的是问题的敏感性，而不是某个求解器已经产生的误差。
\end{proof}

\section{数值等价、消去与求和}

令 $x=10^{16}$。直接计算
\[
  \sqrt{x+1}-\sqrt{x}
\]
会让两个已经舍入为同一浮点数的大数相减，直接相减得到 $0$。乘以共轭式得到
\[
  \sqrt{x+1}-\sqrt{x}
  =\frac{1}{\sqrt{x+1}+\sqrt{x}},
\]
稳定改写与高精度值均为 $5.000\times10^{-9}$。这里精确问题并不病态，失败来自先舍入再相减。

% Generated by tools/build_figures.py; edit figures/figure-specs.json instead.
\MFQFigure{tex/figures/generated/upper-ch14-cancellation.pdf}{两个接近的大数相减会放大输入舍入误差；数学等价的改写可能具有完全不同的数值稳定性。}{upper-ch14-cancellation}


同一原则产生一组常用稳定原语：小 $z$ 时用 \texttt{log1p(z)} 代替 $\log(1+z)$，用 \texttt{expm1(z)} 代替 $e^z-1$；用 \texttt{hypot(x,y)} 避免先平方导致的上溢/下溢。对 $a_i$ 计算
\[
  \log\sum_i e^{a_i}=m+\log\sum_i e^{a_i-m},\qquad m=\max_i a_i,
\]
能把最大指数平移到零。对 $(1000,1000)$，朴素指数上溢，而稳定 log-sum-exp 为 $1000.693147$。稳定改写不是“加大精度”，而是避免丢失本来可保留的信息。

\MFQLead{求和与可复现归约}

对 $(10^{16},1,\ldots,1,-10^{16})$（中间有 10 个 1）按给定顺序计算，朴素、成对、补偿与 Decimal 求和分别得到 $0$、$8$、$10$、$10$。成对求和减少了误差却仍丢失两个单位，说明它不是万能的精确累加器。原因是浮点加法不满足结合律。对 $n$ 项朴素求和，经典界写成
\[
  |\widehat s-s|\le \gamma_{n-1}\sum_i|x_i|,
  \qquad \gamma_k=\frac{ku}{1-ku},
\]
前提是 $ku<1$ 且没有异常值。成对求和把误差增长从线性层数降为对数层数；Kahan/Neumaier 用补偿项保存被舍弃的低位，但不能保证任意输入都正确舍入。

\MFQLead{误差界从哪里来}

设每一步舍入满足 $\operatorname{fl}(a+b)=(a+b)(1+\delta)$、$|\delta|\le u$，并令 $\widehat s_k$ 是前 $k$ 项的左到右和。递推展开得到
\[
 \widehat s_n
 =x_1\prod_{j=2}^{n}(1+\delta_j)
 +\sum_{i=2}^n x_i\prod_{j=i}^{n}(1+\delta_j).
\]
当 $(n-1)u<1$ 时，乘积与 1 的差可由
\[
 \left|\prod_{j=1}^{n-1}(1+\delta_j)-1\right|
 \le\gamma_{n-1}
\]
控制，于是得到上一式的 $\gamma_{n-1}$ 界。证明同时说明两个限制：异常值会破坏单次运算模型，$n$ 太大时 $\gamma_{n-1}$ 不再是小量。数值测试应把输入规模、求和顺序和是否存在上溢分别记录。

\section{线性系统与最小二乘}

对 $Ax=b$，二范数条件数是 $\kappa_2(A)=\|A\|_2\|A^{-1}\|_2=\sigma_{\max}/\sigma_{\min}$。考虑
\[
 A=\begin{pmatrix}1&1\\1&1+10^{-8}\end{pmatrix},
 \qquad x=(1,1)^T,\qquad b=Ax.
\]
由对称矩阵解析特征值可得条件数约为 $4.000\times10^8$。仅把 $b_2$ 增加 $10^{-10}$，精确解变为 $(0.99,1.01)$；前向误差为 $1.000\times10^{-2}$，相对残差不超过 $10^{-15}$。小残差说明计算结果精确满足一个邻近右端项，不说明原问题的答案稳定。

一般线性系统使用带主元的 LU；对对称正定矩阵可用 Cholesky；不要显式形成 $A^{-1}b$。行/列单位差异也会制造数值困难。对
\[
  A_s=\begin{pmatrix}10^{-8}&1\\2\times10^{-8}&3\end{pmatrix}
\]
按列二范数缩放后，列缩放把条件数从 $1.000\times10^9$ 降到 $1.407\times10^1$。缩放改善算法工作坐标，却不改变研究变量的经济单位；报告结果时必须映射回原坐标。

\MFQLead{三种分解各自保存什么}

带主元 LU 在每一步选择当前列中绝对值最大的候选主元并交换行，再消去下方元素；主元过小或矩阵接近奇异时，增长因子和误差都可能变大。Cholesky 只适用于对称正定矩阵，
\[
 A=LL^T,
\]
它利用结构、成本较低，但把不正定矩阵强行送入会在平方根处失败。Householder QR 用正交反射构造 $X=QR$，正交变换不放大二范数；SVD 进一步显示所有奇异方向，适合秩亏与正则化，但计算成本较高。选择分解不是 API 偏好，而是对矩阵结构、可辨识方向和允许误差的声明。

一个实用的后向误差检查是先求候选 $\widehat x$，再计算
\[
 r=b-A\widehat x,\qquad
 \eta_{\rm b}=\frac{\|r\|}{\|A\|\,\|\widehat x\|+\|b\|}.
\]
若 $\eta_{\rm b}$ 很小而前向误差很大，通常是问题病态；若后向误差也大，则算法或实现本身需要修复。这个分解比只打印残差或只打印“求解成功”更能指导下一步。

\MFQLead{最小二乘：不要用正规方程偷走有效数字}

超定系统 $\min_\beta\|X\beta-y\|_2$ 的正规方程为 $X^TX\beta=X^Ty$。若 $X$ 满列秩，
\[
  \kappa_2(X^TX)=\kappa_2(X)^2,
\]
所以先形成 $X^TX$ 会平方条件数。固定实验中，正规方程把条件数从 $2.449\times10^6$ 放大到约 $6.000\times10^{12}$，并产生比 QR 更大的系数误差。QR 直接正交化列空间；SVD
\[
  X=U\Sigma V^T,
  \qquad X^+=V\Sigma^+U^T
\]
能显式显示秩与不可识别方向，代价通常更高。证据程序分别形成正规方程、显式 reduced QR 和显式 SVD 三条路径，而不是把库的 \texttt{lstsq} 实现名称猜成 QR；三条路径均报告系数误差、相对残差与有效秩。

矩阵 $[(1,2),(2,4),(3,6)]$ 的秩为 1；第二个奇异值仅约 $7.32\times10^{-16}$。此时“求得一组系数”不等于因子暴露被识别。可选策略包括删除重复特征、截断小奇异值，或用 ridge 解
\[
  \widehat\beta_\lambda=(X^TX+\lambda I)^{-1}X^Ty.
\]
正则化改变了估计目标，是偏差--方差取舍，不是免费修复。协方差矩阵近奇异时也应先检查谱、样本量和共线暴露，再决定收缩或约束。

\section{容差、测试与研究决策}

常用标量比较为
\[
 |a-b|\le \mathrm{atol}+\mathrm{rtol}\max(|a|,|b|).
\]
$\mathrm{atol}$ 管理接近零的绝对尺度，$\mathrm{rtol}$ 管理远离零的相对尺度。固定实验中，相对策略在 $10^{-8}$ 与 $10^8$ 两个尺度保持一致；纯绝对容差会接受小尺度 100\% 误差，却拒绝大尺度的微小相对误差。数组比较还必须拒绝意外 NaN/Inf，避免 NaN 传播被“都不相等”或特殊选项掩盖。

数值测试不应只比较最终系数。线性求解应检查维度与有限性、相对残差、条件数/奇异值、已知构造解和输入扰动；最小二乘应检查正规方程残差与留出预测；优化应检查可行性、KKT/对偶证书和数据扰动。属性测试可以验证缩放等变性、置换不变性或误差随精度收敛。每个阈值都要记录单位、推导来源与适用尺度。

\begin{example}[残差小而答案错]
令 $A$ 为本章的近奇异矩阵，先用精确 $x=(1,1)^T$ 生成 $b$，再把 $b_2$ 改动 $10^{-10}$。一个稳定算法可能得到残差约 $10^{-15}$，但解从 $(1,1)$ 变为约 $(0.99,1.01)$。若报告只列残差，就会把“精确解了邻近问题”误写成“精确恢复了原参数”；同时列出条件数、扰动方向和前向误差，才能解释这个结果。
\end{example}

\section{输入误差与舍入误差怎样分别观察}
取 $A_\epsilon=\left(\begin{smallmatrix}1&1\\1&1+\epsilon\end{smallmatrix}\right)$。两条方程相减得 $\epsilon x_2=b_2-b_1$；因此 $b_2$ 的扰动 $\delta$ 使 $x_2$ 增加 $\delta/\epsilon$，$x_1$ 减少相同量。这是精确代数关系，即使用无限精度也存在。相反，计算 $\sqrt{x+1}-\sqrt x$ 的消去损失可以用共轭改写显著改善，说明它主要来自算法。

notebook 先用十进制高精度计算这两个例子的参照，再改变 $\epsilon$ 与输入扰动。在结果表中分别列出对原右端项的残差、对扰动右端项的残差与解的变化。若把这三列混为一个“误差”，就无法决定应改算法还是重新考虑参数可识别性。

\noindent\textbf{延伸阅读。}浮点误差、稳定性与病态问题的系统处理参见\cite{higham2002accuracy,golub2013matrix}。

\section{四级问题与即时练习}
\begingroup\small
\begin{enumerate}[nosep,leftmargin=*]
  \item \textbf{口述概念：}区分机器精度、前向/后向误差、问题条件性、算法稳定性和残差。
  \item \textbf{笔试推导：}有理化消去误差公式，并由二阶对称矩阵特征值推导条件数与扰动放大。
  \item \textbf{数值编程：}用 Decimal、补偿求和和解析线性解建立 oracle；覆盖尺度变化、近奇异输入和错误容差。
  \item \textbf{研究判断：}回归或组合优化器报告成功且残差极小，为何结果仍可能不可用？应追加哪些证据？
\end{enumerate}
\endgroup

\subsection*{分级提示}
\begingroup\small
\begin{enumerate}[nosep,leftmargin=*]
  \item \textbf{口述概念：}条件性属于精确问题，稳定性属于算法；残差更接近后向误差。
  \item \textbf{笔试推导：}乘以共轭式；二阶对称矩阵的条件数是最大与最小特征值之比。
  \item \textbf{数值编程：}oracle 不使用同一 float64 路径；容差必须用跨尺度成对案例验证。
  \item \textbf{研究判断：}报告条件数、奇异值、输入扰动、前向误差、残差、正则化路径与样本外稳定性。
\end{enumerate}
\endgroup


\enlargethispage{3\baselineskip}
\noindent\textbf{本章要点。}表示层决定可存储数，稳定性约束算法，条件数约束问题；残差只是证据链的一环。可靠的数值结论应同时给出独立基准、尺度与扰动测试、条件性说明、算法选择和容差来源。
